% GENERATED FILE. Edit canonical lessons, then run npm run generate:books.
% canonical-source-hash: fnv1a64:efeffbc5c17d0540
% canonical-lessons: GU-C462-cancal, GU-C462-tophani, GU-C462-capal, GU-C462-mojilu, GU-C462-vanthvu, GU-R462-review-455-462

\chapter{A Lively Child}
\label{ch:gu-a-lively-child}
\hlchaptermodality{462}

\begin{chapteropening}
\textbf{By the end of this chapter:} \emph{I can describe a lively, mischievous child.}

Complete the last lesson of chapter 462: I can describe a lively, mischievous child.
\end{chapteropening}

\section[\texorpdfstring{\gu{ચંચળ} (cancaḷ)}{ચંચળ (cancaḷ)}]{\gu{ચંચળ} (cancaḷ) — never still; fickle}

\label{lesson:GU-C462-cancal}



\begin{quote}
Before the new one: say the Gujarati for charity, then the Gujarati for a donation.
\end{quote}

\subsection*{You'll want to know: \gu{ચંચળ}}
\textbf{\gu{ચંચળ}} — \emph{cancaḷ} — \textquotedblleft{}never still; fickle\textquotedblright{}.

\begin{quote}
\gu{નાનો} \gu{ભત્રીજો} \gu{ઘણો} \gu{ચંચળ} \gu{છે}.
\end{quote}

\begin{quote}
\emph{nāno bhatrījo ghaṇo cancaḷ che.}
\end{quote}

\begin{quote}
My little nephew is never still.
\end{quote}

\begin{cousinweb}
\textbf{\gu{ચંચળ}} (\emph{cancaḷ}) is a tatsama, Sanskrit \emph{cañcala}, \textquotedblleft{}moving to and fro, unsteady\textquotedblright{}, as Monier-Williams records it.
\end{cousinweb}

\begin{grammarlens}[title={What you've built}]
One more word for describing what you lived and felt.
\end{grammarlens}

\subsection*{Your turn}
\begin{itemize}
\raggedright
  \item \emph{Say it:} \emph{cancaḷ}
  \item \emph{Say it:} \emph{cancaḷ}, once more
  \item \emph{Say it:} \emph{nāno bhatrījo ghaṇo cancaḷ che.} — My little nephew is never still.
  \item \emph{Recall:} say the Gujarati for charity, then the Gujarati for a donation, then say \emph{cancaḷ} again
\end{itemize}

\begin{tcolorbox}[breakable,colback=teal!4,colframe=teal!35!black,title={Before you move on}]
What does \gu{ચંચળ} mean? (\textquotedblleft{}Never still; fickle\textquotedblright{}.) Say it once more.
\end{tcolorbox}
\section[\texorpdfstring{\gu{તોફાની} (tophānī)}{તોફાની (tophānī)}]{\gu{તોફાની} (tophānī) — naughty}

\label{lesson:GU-C462-tophani}



\begin{quote}
Before the new one: say the Gujarati for a donation, then the Gujarati for never still.
\end{quote}

\subsection*{You'll want to know: \gu{તોફાની}}
\textbf{\gu{તોફાની}} — \emph{tophānī} — \textquotedblleft{}naughty\textquotedblright{}.

\begin{quote}
\gu{તોફાની} \gu{છોકરાએ} \gu{મારું} \gu{રમકડું} \gu{તોડ્યું}.
\end{quote}

\begin{quote}
\emph{tophānī chokrāe mārũ ramakḍũ toḍyũ.}
\end{quote}

\begin{quote}
The naughty boy broke my toy.
\end{quote}

\begin{cousinweb}
\textbf{\gu{તોફાની}} (\emph{tophānī}) is built on \textbf{\gu{તોફાન}} (\emph{tophān}), \textquotedblleft{}a storm\textquotedblright{}, traced to Arabic \emph{ṭūfān}, \textquotedblleft{}a flood\textquotedblright{}, through Persian; with \emph{-ī} it names a child who stirs up a storm.
\end{cousinweb}

\begin{grammarlens}[title={What you've built}]
One more word for describing what you lived and felt.
\end{grammarlens}

\subsection*{Your turn}
\begin{itemize}
\raggedright
  \item \emph{Say it:} \emph{tophānī}
  \item \emph{Say it:} \emph{tophānī}, once more
  \item \emph{Say it:} \emph{tophānī chokrāe mārũ ramakḍũ toḍyũ.} — The naughty boy broke my toy.
  \item \emph{Recall:} say the Gujarati for a donation, then the Gujarati for never still, then say \emph{tophānī} again
\end{itemize}

\begin{tcolorbox}[breakable,colback=teal!4,colframe=teal!35!black,title={Before you move on}]
What does \gu{તોફાની} mean? (\textquotedblleft{}Naughty\textquotedblright{}.) Say it once more.
\end{tcolorbox}
\section[\texorpdfstring{\gu{ચપળ} (capaḷ)}{ચપળ (capaḷ)}]{\gu{ચપળ} (capaḷ) — nimble; quick}

\label{lesson:GU-C462-capal}



\begin{quote}
Before the new one: say the Gujarati for never still, then the Gujarati for naughty.
\end{quote}

\subsection*{You'll want to know: \gu{ચપળ}}
\textbf{\gu{ચપળ}} — \emph{capaḷ} — \textquotedblleft{}nimble; quick\textquotedblright{}.

\begin{quote}
\gu{ખિસકોલી} \gu{ઘણી} \gu{ચપળ} \gu{હોય} \gu{છે}.
\end{quote}

\begin{quote}
\emph{khiskolī ghaṇī capaḷ hoy che.}
\end{quote}

\begin{quote}
Squirrels are very nimble.
\end{quote}

\begin{cousinweb}
\textbf{\gu{ચપળ}} (\emph{capaḷ}) is a tatsama, Sanskrit \emph{capala}, \textquotedblleft{}moving, trembling, unsteady\textquotedblright{}, as Monier-Williams records it.
\end{cousinweb}

\begin{grammarlens}[title={What you've built}]
One more word for describing what you lived and felt.
\end{grammarlens}

\subsection*{Your turn}
\begin{itemize}
\raggedright
  \item \emph{Say it:} \emph{capaḷ}
  \item \emph{Say it:} \emph{capaḷ}, once more
  \item \emph{Say it:} \emph{khiskolī ghaṇī capaḷ hoy che.} — Squirrels are very nimble.
  \item \emph{Recall:} say the Gujarati for never still, then the Gujarati for naughty, then say \emph{capaḷ} again
\end{itemize}

\begin{tcolorbox}[breakable,colback=teal!4,colframe=teal!35!black,title={Before you move on}]
What does \gu{ચપળ} mean? (\textquotedblleft{}Nimble; quick\textquotedblright{}.) Say it once more.
\end{tcolorbox}
\section[\texorpdfstring{\gu{મોજીલું} (mojīlũ)}{મોજીલું (mojīlũ)}]{\gu{મોજીલું} (mojīlũ) — fun-loving}

\label{lesson:GU-C462-mojilu}



\begin{quote}
Before the new one: say the Gujarati for naughty, then the Gujarati for nimble.
\end{quote}

\subsection*{You'll want to know: \gu{મોજીલું}}
\textbf{\gu{મોજીલું}} — \emph{mojīlũ} — \textquotedblleft{}fun-loving\textquotedblright{}.

\begin{quote}
\gu{મારો} \gu{મોજીલો} \gu{ભાઈ} \gu{હંમેશા} \gu{ગીતો} \gu{ગાય} \gu{છે}.
\end{quote}

\begin{quote}
\emph{māro mojīlo bhāī hammeśā gīto gāy che.}
\end{quote}

\begin{quote}
My fun-loving brother is always singing.
\end{quote}

\begin{cousinweb}
\textbf{\gu{મોજીલું}} (\emph{mojīlũ}) is built on \emph{moj}, \textquotedblleft{}fun, high spirits\textquotedblright{}, which is usually traced to Arabic \emph{mawj}, \textquotedblleft{}a wave\textquotedblright{}, with the adjective ending \emph{-īlũ}.
\end{cousinweb}

\begin{grammarlens}[title={What you've built}]
One more word for describing what you lived and felt.
\end{grammarlens}

\subsection*{Your turn}
\begin{itemize}
\raggedright
  \item \emph{Say it:} \emph{mojīlũ}
  \item \emph{Say it:} \emph{mojīlũ}, once more
  \item \emph{Say it:} \emph{māro mojīlo bhāī hammeśā gīto gāy che.} — My fun-loving brother is always singing.
  \item \emph{Recall:} say the Gujarati for naughty, then the Gujarati for nimble, then say \emph{mojīlũ} again
\end{itemize}

\begin{tcolorbox}[breakable,colback=teal!4,colframe=teal!35!black,title={Before you move on}]
What does \gu{મોજીલું} mean? (\textquotedblleft{}Fun-loving\textquotedblright{}.) Say it once more.
\end{tcolorbox}
\section[\texorpdfstring{\gu{વંઠવું} (vanṭhvũ)}{વંઠવું (vanṭhvũ)}]{\gu{વંઠવું} (vanṭhvũ) — to go astray; to be spoilt}

\label{lesson:GU-C462-vanthvu}



\begin{quote}
Before the new one: say the Gujarati for nimble, then the Gujarati for fun-loving.
\end{quote}

\subsection*{You'll want to know: \gu{વંઠવું}}
\textbf{\gu{વંઠવું}} — \emph{vanṭhvũ} — \textquotedblleft{}to go astray; to be spoilt\textquotedblright{}.

\begin{quote}
\gu{ખરાબ} \gu{દોસ્તોને} \gu{લીધે} \gu{દીકરો} \gu{વંઠી} \gu{ગયો}.
\end{quote}

\begin{quote}
\emph{kharāb dostone līdhe dīkro vanṭhī gayo.}
\end{quote}

\begin{quote}
Because of bad friends the son went astray.
\end{quote}

Past (he …): \gu{તે} \gu{વંઠી} \gu{ગયો}, \emph{te vanṭhī gayo}.

\begin{cousinweb}
\textbf{\gu{વંઠવું}} (\emph{vanṭhvũ}) is inherited: CDIAL lists the Gujarati verb, \textquotedblleft{}to be dissipated\textquotedblright{}, under Sanskrit \emph{vaṇṭa}, \textquotedblleft{}tailless, unmarried\textquotedblright{} (CDIAL 11236).
\end{cousinweb}

\begin{grammarlens}[title={What you've built}]
One more word for describing what you lived and felt.
\end{grammarlens}

\subsection*{Your turn}
\begin{itemize}
\raggedright
  \item \emph{Say it:} \emph{vanṭhvũ}
  \item \emph{Say it:} \emph{vanṭhvũ}, once more
  \item \emph{Say it:} \emph{kharāb dostone līdhe dīkro vanṭhī gayo.} — Because of bad friends the son went astray.
  \item \emph{Recall:} say the Gujarati for nimble, then the Gujarati for fun-loving, then say \emph{vanṭhvũ} again
\end{itemize}

\begin{tcolorbox}[breakable,colback=teal!4,colframe=teal!35!black,title={Before you move on}]
What does \gu{વંઠવું} mean? (\textquotedblleft{}To go astray; to be spoilt\textquotedblright{}.) Say it once more.
\end{tcolorbox}
\section[(dialogue)]{Looking back: from A Heavy Heart to A Lively Child}

\label{lesson:GU-R462-review-455-462}



\begin{quote}
Say the Gujarati for gloomy and anguish.
\end{quote}

\subsection*{Your turn}
Say these aloud:

\begin{itemize}
\raggedright
  \item gloomy, anguish, a sigh, to make someone cry, weeping
  \item to be hurt, unwelcome, a taunt, disregard, a knock
  \item hatred, to despise, contempt, envy, dislike
  \item the pain of separation, solitude, deserted, to crave, to come to mind
  \item tender, harsh, cruel, good-hearted, to soften with pity
  \item trusting, sly, shrewd, deceit, to cajole
  \item miserly, wasteful with money, penniless, charity, a donation
  \item never still, naughty, nimble, fun-loving, to go astray
\end{itemize}

\begin{tcolorbox}[breakable,colback=teal!4,colframe=teal!35!black,title={Before you move on}]
Gloomy? (\textbf{\gu{ગમગીન}}, \emph{gamgīn}.) Anguish? (\textbf{\gu{વ્યથા}}, \emph{vyathā}.) A sigh? (\textbf{\gu{નિસાસો}}, \emph{nisāso}.) To make someone cry? (\textbf{\gu{રડાવવું}}, \emph{raḍāvvũ}.) Weeping? (\textbf{\gu{રુદન}}, \emph{rudan}.) To be hurt? (\textbf{\gu{દુભાવું}}, \emph{dubhāvũ}.) Unwelcome? (\textbf{\gu{માઠું}}, \emph{māṭhũ}.) A taunt? (\textbf{\gu{મહેણું}}, \emph{maheṇũ}.) Disregard? (\textbf{\gu{અવગણના}}, \emph{avgaṇnā}.) A knock? (\textbf{\gu{ઠેસ}}, \emph{ṭhes}.) Hatred? (\textbf{\gu{નફરત}}, \emph{naphrat}.) To despise? (\textbf{\gu{ધિક્કારવું}}, \emph{dhikkārvũ}.) Contempt? (\textbf{\gu{તિરસ્કાર}}, \emph{tiraskār}.) Envy? (\textbf{\gu{અદેખાઈ}}, \emph{adekhāī}.) Dislike? (\textbf{\gu{અણગમો}}, \emph{aṇgamo}.) The pain of separation? (\textbf{\gu{વિરહ}}, \emph{virah}.) Solitude? (\textbf{\gu{એકાંત}}, \emph{ekā̃t}.) Deserted? (\textbf{\gu{સૂનું}}, \emph{sūnũ}.) To crave? (\textbf{\gu{તરસવું}}, \emph{tarasvũ}.) To come to mind? (\textbf{\gu{સાંભરવું}}, \emph{sā̃bharvũ}.) Tender? (\textbf{\gu{કોમળ}}, \emph{komaḷ}.) Harsh? (\textbf{\gu{કઠોર}}, \emph{kaṭhor}.) Cruel? (\textbf{\gu{ક્રૂર}}, \emph{krūr}.) Good-hearted? (\textbf{\gu{ભલું}}, \emph{bhalũ}.) To soften with pity? (\textbf{\gu{પીગળવું}}, \emph{pīgaḷvũ}.) Trusting? (\textbf{\gu{ભોળું}}, \emph{bhoḷũ}.) Sly? (\textbf{\gu{લુચ્ચું}}, \emph{luccũ}.) Shrewd? (\textbf{\gu{ચાલાક}}, \emph{cālāk}.) Deceit? (\textbf{\gu{કપટ}}, \emph{kapaṭ}.) To cajole? (\textbf{\gu{ફોસલાવવું}}, \emph{phoslāvvũ}.) Miserly? (\textbf{\gu{કૃપણ}}, \emph{kṛpaṇ}.) Wasteful with money? (\textbf{\gu{ઉડાઉ}}, \emph{uḍāu}.) Penniless? (\textbf{\gu{કંગાળ}}, \emph{kangāḷ}.) Charity? (\textbf{\gu{સખાવત}}, \emph{sakhāvat}.) A donation? (\textbf{\gu{દાન}}, \emph{dān}.) Never still? (\textbf{\gu{ચંચળ}}, \emph{cancaḷ}.) Naughty? (\textbf{\gu{તોફાની}}, \emph{tophānī}.) Nimble? (\textbf{\gu{ચપળ}}, \emph{capaḷ}.) Fun-loving? (\textbf{\gu{મોજીલું}}, \emph{mojīlũ}.) To go astray? (\textbf{\gu{વંઠવું}}, \emph{vanṭhvũ}.)
\end{tcolorbox}
